\chapter{Automated Alert Response with Falco Handler}

Detecting an attack in real time is only half the battle. The other half is \textbf{responding} to it. Falco, by itself, is a pure detection engine -- it generates alerts but does not take any automatic action on the cluster. In production environments, security teams integrate Falco with an \textbf{automated response system} to minimize the window between detection and remediation.

In this chapter, we will extend our Falco deployment with two additional components:
\begin{itemize}
    \item \textbf{Falcosidekick}: a companion service that receives alerts from Falco via an internal channel and fans them out to external destinations such as webhooks, Slack, PagerDuty, or S3.
    \item \textbf{Falco Handler}: a custom Python microservice deployed in the cluster that receives Falco alerts via webhook and executes an automated response -- in our case, labeling the compromised deployment with \texttt{tainted-by-falco=true}.
\end{itemize}

\begin{tcolorbox}[colback=blue!5!white,colframe=blue!75!black,title=Detection vs. Response]
The pipeline we are building implements the two phases of an Incident Response cycle:
\begin{itemize}
    \item \textbf{Detection (Falco)}: observes kernel-level events and generates structured alerts.
    \item \textbf{Response (Handler)}: consumes those alerts and triggers an action. The action in this lab is labeling the deployment, but the same architecture can be used to isolate a pod by scaling it to zero replicas, rotate credentials, or trigger a Slack notification.
\end{itemize}
\end{tcolorbox}

\section{Architecture Overview}

The full alert pipeline works as follows:

\begin{enumerate}
    \item Falco intercepts a suspicious syscall and generates a JSON alert.
    \item The alert is sent over an internal HTTP channel to \textbf{Falcosidekick}, which runs as a sidecar container in the same Falco pod.
    \item Falcosidekick forwards the alert as an HTTP POST request to the \textbf{Falco Handler} service endpoint.
    \item The Handler extracts the pod name and namespace from the JSON payload, resolves the parent deployment, and labels it.
\end{enumerate}

\section{Deploying the Falco Handler}

\subsection{The Handler Microservice}

The handler is a minimal Python Flask application. It exposes a single \texttt{POST /} endpoint that receives the JSON alert from Falcosidekick, identifies the affected Kubernetes deployment, and applies a label to it.

The handler requires \textbf{RBAC permissions} to interact with the Kubernetes API. We must create a dedicated ServiceAccount, a ClusterRole granting read access to pods and patch access to deployments, and bind them together.

\begin{tcolorbox}[colback=blue!5!white,colframe=blue!75!black,title=Info]
The handler image and all necessary Kubernetes manifests (\texttt{falco-handler-rbac.yaml}, \texttt{falco-handler-deployment.yaml}, \texttt{falco-handler-service.yaml}) are included in the \texttt{material/falco\_handler} directory.
\end{tcolorbox}

\subsection{Deploying the Handler}

Navigate to the handler's manifests directory and apply the resources in the \texttt{default} namespace:

\begin{lstlisting}[language=bash]
cd ../falco_handler
kubectl apply -f falco-handler-rbac.yaml
kubectl apply -f falco-handler-deployment.yaml
kubectl apply -f falco-handler-service.yaml
\end{lstlisting}

Verify that the handler pod is running:

\begin{lstlisting}[language=bash]
kubectl get pods -l app=falco-handler
\end{lstlisting}

\begin{lstlisting}[language=bash]
NAME                             READY   STATUS    RESTARTS   AGE
falco-handler-7d6f9c8b4-xk2mp   1/1     Running   0          15s
\end{lstlisting}

\subsection{Reconfiguring Falco with Falcosidekick}

We now need to upgrade our Falco Helm release to enable Falcosidekick and point it to the handler's service endpoint. The handler is exposed as a ClusterIP service named \texttt{falco-handler} on port 80.

Create a file \texttt{falco-sidekick-values.yaml}:

\begin{lstlisting}[language={}]
cat <<EOF > falco-sidekick-values.yaml
falcosidekick:
  enabled: true
  config:
    webhook:
      enabled: true
      address: http://falco-handler.default.svc.cluster.local:80/
EOF
\end{lstlisting}

Upgrade the Falco Helm release to apply these settings:

\begin{lstlisting}[language=bash]
helm upgrade falco falcosecurity/falco \
  -n falco \
  --reuse-values \
  -f falco-sidekick-values.yaml
\end{lstlisting}

Wait for the pods to restart:

\begin{lstlisting}[language=bash]
kubectl rollout status daemonset/falco -n falco
\end{lstlisting}

After the rollout, each Falco pod will now have two containers: \texttt{falco} and \texttt{falcosidekick}.

\begin{lstlisting}[language=bash]
kubectl get pods -n falco
\end{lstlisting}

\begin{lstlisting}[language=bash]
NAME          READY   STATUS    RESTARTS   AGE
falco-4dgzk   2/2     Running   0          45s
\end{lstlisting}

\section{Observing the Automated Response}

Trigger any of the attacks from Chapter~3 (e.g., upload the \texttt{ptrace\_payload.zip} or run the reverse shell). After the attack, check the handler logs:

\begin{lstlisting}[language=bash]
kubectl logs deployment/falco-handler
\end{lstlisting}

You should see output similar to the following, showing that the handler received the alert, resolved the deployment name, and applied the label:

\begin{lstlisting}[language=bash]
[*] Triggered rule from pod=zipapp-5ffd444d74-jbwq2, ns=default
[*] Deployment: zipapp
[+] Labeled deployment zipapp in default
\end{lstlisting}

Now verify that the label was applied to the deployment:

\begin{lstlisting}[language=bash]
kubectl get deployments --show-labels
\end{lstlisting}

\begin{lstlisting}[language=bash]
NAME     READY   UP-TO-DATE   AVAILABLE   LABELS
zipapp   1/1     1            1           app=zipapp,tainted-by-falco=true
\end{lstlisting}

The label \texttt{tainted-by-falco=true} marks the deployment as compromised. This label can be used by other systems to trigger further remediation actions.

\begin{tcolorbox}[colback=green!5!white,colframe=green!75!black,title=Question]
\begin{itemize}
    \item The current handler only applies a label. Look at the commented-out code block in the handler: what more aggressive action was planned but left disabled? When would it be appropriate to enable it?
    \item The handler uses a \texttt{ClusterRole} (cluster-wide permissions) rather than a \texttt{Role} (namespace-scoped). Why is this necessary given how the handler resolves the pod's deployment?
    \item After the attack, remove the label from the deployment manually:
\end{itemize}
\begin{lstlisting}[language=bash]
kubectl label deployment zipapp tainted-by-falco-
\end{lstlisting}
\end{tcolorbox}

